\chapter{Opioid Withdrawal Symptoms}

\medskip
\noindent\textit{\textbf{Under review.} This chapter is an AI-assisted educational draft; it is not a substitute for professional medical advice.}

\section*{Definition and Overview}
Opioid withdrawal is a predictable group of symptoms that can occur when a person dependent on opioids suddenly reduces or stops them. It is usually very uncomfortable and is rarely fatal by itself, but dehydration, relapse, overdose after loss of tolerance, and coexisting illness can be dangerous.

\section*{Clinical Features}
Symptoms may include craving, anxiety, restlessness, muscle aches, sweating, runny nose, yawning, abdominal cramps, nausea, vomiting, diarrhoea, dilated pupils, gooseflesh, insomnia, and rapid heartbeat. They vary with the opioid, dose, duration, and time since the last dose.

\section*{When to Seek Medical Care}
Seek urgent help for severe vomiting or diarrhoea, dehydration, chest pain, confusion, pregnancy, severe mental distress, or suspected overdose. Call 000 for slow or stopped breathing, blue lips, collapse, seizure, or unconsciousness.

\section*{Assessment and Management}
A clinician assesses opioid use, other substances, pregnancy, mental health, pain, and medical complications. Treatment may include supervised opioid agonist therapy such as buprenorphine or methadone, symptom medicines, fluids, and psychosocial support. Do not self-manage with another person's medicine.

\section*{Prevention and Self-Care}
Do not restart a previous opioid dose after withdrawal without medical advice because tolerance falls and overdose risk rises. Keep naloxone available where appropriate, avoid alcohol and sedatives, and connect with an alcohol and other drug service for ongoing support.

\section*{Expanded clinical context}
Opioid withdrawal symptoms is a mental-health, behavioural, or support topic. Interpretation should account for age, baseline health, medicines, comorbidities, goals, and access to care. This chapter supports discussion with a qualified clinician and does not establish a diagnosis.

\section*{Assessment and decision points}
Review onset, duration, triggers, sleep, substance use, physical health, developmental context, social stressors, and immediate safety. Ask about self-harm or harm from others when appropriate. Document the main concern, time course, functional impact, relevant risks, and red flags. Use targeted investigations when their result is likely to change management.

\section*{Management and follow-up}
Care may include education, psychological therapy, practical support, treatment of coexisting conditions, medication when indicated, and a shared follow-up plan. Explain expected improvement, when reassessment is needed, and how follow-up can be accessed. Avoid starting, stopping, or sharing prescription medicines without professional advice.

\section*{Safety-netting}
Seek urgent help for immediate danger, suicidal intent, psychosis, severe confusion, inability to meet basic needs, or rapidly worsening symptoms.

\section*{Practical prevention}
Use plain-language information, supportive relationships, regular routines, and professional review when symptoms persist or impair daily life. Keep written instructions and seek review if the topic recurs, changes, or does not improve as expected.

\section*{Limitations}
This educational expansion should be checked against current local guidance and authoritative topic-specific sources.
\section*{Communication and support}
Withdrawal is easier to manage when the person has a clear plan, trusted support, and rapid access to treatment if symptoms or cravings escalate. Involve the person in decisions and record emergency contacts, follow-up arrangements, and overdose-prevention advice.
\clearpage
